\documentclass{article}

\usepackage[utf8]{inputenc}
\usepackage[T1]{fontenc}
\usepackage{lmodern}
\usepackage[margin=1in]{geometry}
\usepackage{hyperref}
\usepackage{amsmath,amssymb,amsfonts}
\usepackage{booktabs}
\usepackage{microtype}

\newcommand{\MPL}{\mathrm{MPL}}
\newcommand{\LD}{\mathrm{LD}}
\newcommand{\obs}{\mathrm{obs}}
\newcommand{\cosine}{\mathrm{cos}}

\hypersetup{
  pdftitle={Supplementary Note: q2 Half-Life MPL-LD Response},
  pdfauthor={Jiaju Wu}
}

\title{Supplementary Note:\\
The \(q_2\) Half-Life MPL-LD Residual Response}
\author{Jiaju Wu}
\date{June 19, 2026}

\begin{document}
\maketitle

\section{Purpose}
This note records the current theoretical interpretation of the residual correction used in
the main paper.  It replaces the older DropRelaxS-style supplement.  The goal is not to
prove a universal optimizer law, but to make explicit why the current estimator has the form
\begin{equation}\label{eq:supp_main}
  \widehat L_s(t)
  =
  L_{\MPL,s}(t)
  +
  a_s\widehat\kappa_s\phi_{\lambda_s,s}(t).
\end{equation}
The three structural choices are:
\begin{enumerate}
  \item a causal LR-drop response feature \(\phi_{\lambda_s,s}\);
  \item a schedule-only \(q_2\) half-life rule for \(\lambda_s\);
  \item MPL-LD tangent projection before estimating the single scalar \(\widehat\kappa_s\).
\end{enumerate}

\section{Linear Response Approximation}
Let \(e_t\) be an unobserved excess-loss tracking state after the MPL baseline has modeled the
dominant quasi-static trend.  Around a fixed source schedule, a first-order local model is
\begin{equation}\label{eq:state}
  e_t=A_t e_{t-1}+b\,d_t+\epsilon_t,
  \qquad
  d_t=[\eta_{t-1}-\eta_t]_+ .
\end{equation}
The positive LR decrement \(d_t\) is the forcing term: when the LR drops, the ideal
LR-dependent loss floor moves, while the observed training state responds with finite memory.
For a stable one-pole approximation in cumulative-LR time,
\begin{equation}
  A_t\approx \exp(-\lambda_s\eta_t).
\end{equation}
Substituting this into Eq.~\eqref{eq:state} gives
\begin{equation}\label{eq:z}
  z_t=\exp(-\lambda_s\eta_t)z_{t-1}+d_t,\qquad
  \phi_{\lambda_s,s}(t)=z_t/\eta_{\max}.
\end{equation}
Equivalently,
\begin{equation}
  \phi_{\lambda_s,s}(t)
  =
  \frac{1}{\eta_{\max}}
  \sum_{u\le t}d_u
  \exp\!\left(-\lambda_s\sum_{v=u+1}^{t}\eta_v\right).
\end{equation}
This derivation explains the qualitative shape of the feature: it is causal, drop-only, and
measured in cumulative LR.  The scalar \(\kappa\) absorbs the unknown mapping from the hidden
state to loss units.

\section[q2 Half-Life Rule]{\(q_2\) Half-Life Rule}
The response rate is not fitted from WSD targets.  It is computed from the target LR schedule.
For \(D_s=\sum_t d_t\), define
\begin{equation}
  p_t=\frac{d_t}{D_s},\qquad
  q_s=\sum_t p_t^2=\frac{1}{n_{\mathrm{eff}}}.
\end{equation}
Thus \(q_s\) is the inverse effective number of drop atoms.  A single sharp drop has
\(q_s=1\); a drop uniformly spread over \(n\) steps has \(q_s=1/n\).

The observed curves are logged mainly every \(\Delta_{\obs}=128\) steps.  We set
\begin{equation}
  \lambda_{\obs}
  =
  \frac{\log2}{\eta_{\max}\Delta_{\obs}},
\end{equation}
and use the half-life bracket
\begin{equation}\label{eq:halflife}
  H_s=(2-q_s)\Delta_{\obs},
  \qquad
  \lambda_s=\frac{\log2}{\eta_{\max}H_s}
  =
  \frac{\lambda_{\obs}}{2-q_s}.
\end{equation}
The interpretation is simple: diffuse decay is assigned roughly a two-observation half-life,
while a single-step drop is assigned a one-observation half-life.  This is a resolution-based
structural prior, not a fitted physical constant.

\section{Support-Projection Locality}
The correction should represent local cooldown forcing, not full-horizon diffuse decay.  Let
the post-warmup horizon be \(H=T_s-W\), and let \(\ell_s\) be the positive-drop support span.
Let \(m_s\) be the uniform density on that support, and let \(u\) be the uniform density on
the full horizon.  Projecting away the diffuse mode gives
\begin{equation}
  a_s
  =
  \mathbf 1\{D_s>0\}
  \frac{\|(I-P_u)m_s\|_2^2}{\|m_s\|_2^2}.
\end{equation}
Since \(\|m_s\|_2^2=1/\ell_s\) and \(\|P_um_s\|_2^2=1/H\),
\begin{equation}\label{eq:locality_supp}
  a_s
  =
  \mathbf 1\{D_s>0\}
  \left[1-\frac{\ell_s}{T_s-W}\right]_+ .
\end{equation}
This is why the term is not a learned gate.  It is an explicit energy fraction after removing
the global diffuse component.

\section{MPL-LD Nuisance Projection}
The source residual is
\begin{equation}
  r_{\cosine}(t)=L_{\cosine}(t)-L_{\MPL,\cosine}(t).
\end{equation}
If the LR-dependent MPL parameters are locally perturbed, the residual contains
\begin{equation}
  J_{\LD}(t)\Delta\theta_{\LD},
  \qquad
  J_{\LD}(t)=
  \left[
  \frac{\partial L_{\MPL}}{\partial\log B},
  \frac{\partial L_{\MPL}}{\partial\log C},
  \frac{\partial L_{\MPL}}{\partial\log\beta},
  \frac{\partial L_{\MPL}}{\partial\log\gamma}
  \right].
\end{equation}
Let \(P_{\LD}\) be the orthogonal projection onto this four-dimensional tangent space on the
cosine calibration suffix.  For target schedule \(s\),
\begin{equation}
  x_s=(I-P_{\LD})\phi_{\lambda_s,\cosine},
  \qquad
  y=(I-P_{\LD})r_{\cosine}.
\end{equation}
The fitted amplitude is
\begin{equation}\label{eq:kappa_supp}
  \widehat\kappa_s
  =
  \frac{\langle x_s,y\rangle_+}
       {\|x_s\|_2^2+1/N_{\mathrm{cal}}}.
\end{equation}
Only this scalar is estimated from loss residuals.  The response rate, locality factor, and
nuisance subspace are all determined by the LR schedule and frozen MPL formula.

\section{Current Evidence and Boundaries}
The current theory-refinement audit reports:
\begin{center}
\begin{tabular}{lrrr}
\toprule
split & mean MAE change & worst row & wins \\
\midrule
same-scale WSD-family & \(-29.88\%\) & \(-4.67\%\) & 15/15 \\
cross-scale WSD-family & \(-24.95\%\) & \(-3.15\%\) & 30/30 \\
controls & \(+0.00\%\) & \(+0.00\%\) & 9/9 non-harm \\
\bottomrule
\end{tabular}
\end{center}
The main negative controls are equally important.  Without MPL-LD nuisance projection, the raw
one-dimensional projection fails on WSD-family targets.  Without support-projection locality,
the WSD mean improves slightly but the short-cosine control is harmed.  Density projection is
natural but still keeps too much diffuse cosine signal.  These failures justify the current
low-capacity formula and also define its limitations.

\end{document}
